

\label{02-chapter}

In questo lavoro di tesi verranno analizzati tutti i vari passaggi necessari per la realizzazione di un guanto capace di rivelare quando l'indossatore schiocca le dita.
Lo scopo è quello di fornire le basi per la realizzazione di un generico algoritmo di Tiny Machine Learning (TinyML), che potrà poi essere portato su un qualsiasi microcontrollore.\\

\noindent{}Ho deciso di dividere la tesi in due capitoli principali: uno per la descrizione della strumentazione utilizzata, e l'altro dedicato alla spiegazione di come ho realizzato il progetto.\\

\noindent{}Nel primo capitolo saranno elencati tutti i dispositivi utilizzati, soffermandosi anche sulla teoria dietro al funzionamento dei sensori con cui vengono raccolti i dati.\\

\noindent{}Nel secondo capitolo verrà spiegato invece come acquisire i dati dai sensori, creare ed addestrare una rete neurale, ed infine caricarla su una scheda Arduino. Verranno inoltre riportate le fondamentali funzioni Python ed Arduino utilizzate nei vari codici, e la loro implementazione nel progetto.\\